% !TEX root = paper.tex
\section{Deferred PIOPs}
\label{sec:deferred-piops}
This appendix collects the low-level PIOPs and relations used throughout \papername{}. These are standard or lightly adapted from prior work, and are deferred here so that the main text can focus on the string-filtering protocols.

\subsection{NoZero Check}
The \emph{NoZero Check} takes a column $\col = (\mathsf{d},\mathsf{a})$ and proves that every active entry of $\col$ is nonzero. We package this as the relation
\begin{displaymath}
\NPRelation_{\mathsf{NoZero}} = \left\{\,
  \NPInstance = n,\;
  \NPInstancey = (\Oracle{\MLE{\col}}),\;
  \NPWitness = (\col)
  \;\middle|\;
    \forall i,\, \mathsf{a}[i] = 1 \implies \col[i] \neq 0
\,\right\}.
\end{displaymath}

A field element is nonzero exactly when it has a multiplicative inverse. The prover sends an inverse column $\mathsf{inv}$ whose value at each active index is the inverse of the data there and which is $0$ at inactive indices. If $\mathsf{inv}$ is formed this way, then $\mathsf{d}\mult\mathsf{inv} = \mathsf{a}$ holds across the whole hypercube: at an active index both sides equal $1$, and at an inactive index both sides equal $0$ (since $\mathsf{inv}$ is $0$ there). A single \cref{piop:zerocheck} on $\mathsf{d}\mult\mathsf{inv} - \mathsf{a}$ certifies this, and forces every active entry to be invertible, hence nonzero.
\begin{DefinePIOP}{NoZero Check}{nozero-check}
  \emph{Inputs:} $\Oracle{\MLE{\col}} = (\Oracle{\MLE{\mathsf{a}}}, \Oracle{\MLE{\mathsf{d}}})$.
  \begin{enumerate}[nosep]
    \item $\PIOPProver$ constructs and sends the inverse oracle $\Oracle{\MLE{\mathsf{inv}}}$, with $\mathsf{inv}[i] = \mathsf{d}[i]^{-1}$ at active indices and $\mathsf{inv}[i] = 0$ elsewhere.
    \item $\PIOPProver$ and $\PIOPVerifier$ invoke \cref{piop:zerocheck} on the column $\mathsf{d}\mult\mathsf{inv} - \mathsf{a}$.
  \end{enumerate}
\end{DefinePIOP}

\subsection{Booleanity Check}
The \emph{Booleanity Check} takes a column $\col = (\mathsf{d},\mathsf{a})$ and proves that every active entry of $\col$ is $0$ or $1$. We package this as the relation
\begin{displaymath}
\NPRelation_{\mathsf{Bool}} = \left\{\,
  \NPInstance = \bot,\;
  \NPInstancey = (\Oracle{\MLE{\col}}),\;
  \NPWitness = (\col)
  \;\middle|\;
    \forall i,\, \mathsf{a}[i] = 1 \implies \mathsf{d}[i]\in\{0,1\}
\,\right\}.
\end{displaymath}
Since $\mathsf{d}[i]\in\{0,1\}$ exactly when $\mathsf{d}[i]\bigl(1-\mathsf{d}[i]\bigr)=0$, the check is a single \cref{piop:zerocheck} on the column $\mathsf{a}\cdot\mathsf{d}\cdot(1-\mathsf{d})$, which vanishes everywhere iff every active entry is boolean.
\begin{DefinePIOP}{Booleanity Check}{boolean-check}
  \emph{Inputs:} $\Oracle{\MLE{\col}} = (\Oracle{\MLE{\mathsf{a}}}, \Oracle{\MLE{\mathsf{d}}})$.
  \begin{enumerate}[nosep]
    \item $\PIOPProver$ and $\PIOPVerifier$ invoke \cref{piop:zerocheck} on the column $\mathsf{a}\cdot\mathsf{d}\cdot(1-\mathsf{d})$.
  \end{enumerate}
\end{DefinePIOP}

\subsection{No-Duplicate Check}
The \emph{No-Duplicate Check} takes a column $\col = (\mathsf{d},\mathsf{a})$ and proves that the active entries of $\col$ are pairwise distinct. We package this as the relation
\begin{displaymath}
\NPRelation_{\mathsf{NoDup}} = \left\{\,
  \NPInstance = \bot,\;
  \NPInstancey = (\Oracle{\MLE{\col}}),\;
  \NPWitness = (\col)
  \;\middle|\;
    \forall i\neq j,\; \mathsf{a}[i]=\mathsf{a}[j]=1 \implies \mathsf{d}[i]\neq\mathsf{d}[j]
\,\right\}.
\end{displaymath}

\begin{DefinePIOP}{No-Duplicate Check}{nodup-check}
    TODO
\end{DefinePIOP}

\subsection{Limb Reconstruction Check}
The \emph{Limb Reconstruction Check} takes 1 $n$-bit column $\col$ and $k$ columns $\col_0,\dots,\col_{k-1}$ each with $n/k$ bits and proves that each entry of $\col$ is the concatenation of the corresponding entries in $\col_0,\dots,\col_{k-1}$. We package this as the relation
\begin{displaymath}
\NPRelation_{\mathsf{Limb}} = \left\{\,
  \NPInstance = n,\;
  \NPInstancey = (\Oracle{\MLE{\col}}, \Oracle{\MLE{\col_0}}, \dots, \Oracle{\MLE{\col_{k-1}}}),\;
  \NPWitness = (\col, \col_0, \dots, \col_{k-1})
  \;\middle|\;
    \forall i,\, \col[i] = \sum_{j=0}^{k-1} 2^{(k-1-j)n/k} \col_j[i]
\,\right\}.
\end{displaymath}


Reconstruction is a single linear identity over the hypercube: subtracting the weighted limbs from $\col$ must vanish at every row. A single \cref{piop:zerocheck} on $\col - \sum_{j=0}^{k-1} 2^{(k-1-j)n/k}\mult\col_j$ enforces exactly this, pinning down each entry of $\col$ as the concatenation of the corresponding limbs. The check certifies only the weighted-sum relationship; bounding each limb $\col_j$ to its $n/k$-bit width, when required, is done by a separate \cref{piop:range-check}.
\begin{DefinePIOP}{Limb Reconstruction Check}{limb-check}
  \emph{Inputs:} $\Oracle{\MLE{\col}}, \Oracle{\MLE{\col_0}}, \ldots, \Oracle{\MLE{\col_{k-1}}}$.
  \begin{enumerate}[nosep]
    \item $\PIOPProver$ and $\PIOPVerifier$ invoke \cref{piop:zerocheck} on the column $\col - \sum_{j=0}^{k-1} 2^{(k-1-j)n/k}\mult\col_j$.
  \end{enumerate}
\end{DefinePIOP}

\subsection{Lookup Check}
The \emph{Lookup Check} takes two columns $\col_1 = (\mathsf{d}_1, \mathsf{a}_1)$ and $\col_2 = (\mathsf{d}_2, \mathsf{a}_2)$ and proves that every active value of $\col_1$ also occurs as an active value of $\col_2$, written $\col_1 \lookup \col_2$. We package this as the relation
\begin{displaymath}
\NPRelation_{\mathsf{Lookup}} = \left\{\,
  \NPInstance = \bot,\;
  \NPInstancey = (\Oracle{\MLE{\col_1}}, \Oracle{\MLE{\col_2}}),\;
  \NPWitness = (\col_1, \col_2)
  \;\middle|\;
    \forall i,\, \mathsf{a}_1[i] = 1 \implies \exists j:\, \mathsf{a}_2[j] = 1 \wedge \mathsf{d}_1[i] = \mathsf{d}_2[j]
\,\right\}.
\end{displaymath}
We discharge the lookup with the LogUp technique~\cite{Habock22}: $\col_1 \lookup \col_2$ holds if and only if there is a \emph{multiplicity} column $\mathsf{m}$ over $\col_2$'s hypercube such that, with $\mu$ and $\nu$ the hypercube dimensions of $\MLE{\col_1}$ and $\MLE{\col_2}$,
\begin{equation}
  \label{eq:lookup-logup}
  \sum_{x \in \BoolHCube{\mu}} \frac{\MLE{\mathsf{a}_1}(x)}{X - \MLE{\mathsf{d}_1}(x)}
  \;=\;
  \sum_{y \in \BoolHCube{\nu}} \frac{\MLE{\mathsf{a}_2}(y)\mult\MLE{\mathsf{m}}(y)}{X - \MLE{\mathsf{d}_2}(y)}
\end{equation}
as rational functions in $X$, where the activators zero out the inactive terms on each side. The prover sends $\mathsf{m}$, and the parties verify the rational identity by Schwartz--Zippel: the verifier samples a random $r \samples \F$, the prover sends the purported common value $s$ that both sides take at $X = r$, and the parties run \cref{piop:ratsumcheck} twice to confirm that each side evaluates to $s$ at $X = r$. When $\col_2$ is known to be duplicate-free, $\mathsf{m}(y)$ records the multiplicity of $\MLE{\mathsf{d}_2}(y)$ in $\col_1$, so the verifier stores $\Oracle{\MLE{\mathsf{m}}}$ as a protocol output for use by higher-level checks.
\begin{DefinePIOP}{Lookup Check}{lookup-check}
  \emph{Inputs:} $\Oracle{\MLE{\col_1}}, \Oracle{\MLE{\col_2}}$.
  \begin{enumerate}[nosep]
    \item $\PIOPProver$ constructs and sends the multiplicity oracle $\Oracle{\MLE{\mathsf{m}}}$, which $\PIOPVerifier$ stores as protocol output.
    \item $\PIOPVerifier$ samples a random point $r \samples \F$ and sends it to $\PIOPProver$, who responds with the claimed common value $s$ of both sides of \cref{eq:lookup-logup} at $X = r$.
    \item $\PIOPProver$ and $\PIOPVerifier$ invoke \cref{piop:ratsumcheck} twice to prove that both sides of \cref{eq:lookup-logup} evaluate to $s$ at $X = r$.
  \end{enumerate}
\end{DefinePIOP}


\subsection{Range Check}
The \emph{Range Check} takes a column $\col = (\mathsf{d}, \mathsf{a})$ and a width $n$, and proves that every active entry of $\col$ is a non-negative integer of bitwidth $n$, i.e.\ lies in $[0, 2^n - 1]$. We package this as the relation
\begin{displaymath}
\NPRelation_{\mathsf{Range}} = \left\{\,
  \NPInstance = n,\;
  \NPInstancey = (\Oracle{\MLE{\col}}),\;
  \NPWitness = (\col)
  \;\middle|\;
    \forall i,\, \mathsf{a}[i] = 1 \implies 0\leq \mathsf{d}[i] < 2^n
\,\right\}.
\end{displaymath}
We assume $n < \log_2 |\F|$, so the range embeds unambiguously into $\F$. A direct approach proves $\col \lookup \mathsf{rng}_n := (\mathbf{1}, \idx_n)$, since the index polynomial $\idx_n$ ranges over exactly $\{0, \ldots, 2^n - 1\}$; but this lookup table has $2^n$ rows and is wasteful for large $n$. Instead we split each value into $16$-bit \emph{limbs} and range-check each limb against a small index column. Choosing $k$ and $\beta$ with $n = 16k + \beta$ and $\beta \leq 16$, the prover sends limb columns $\mathsf{d}_1, \ldots, \mathsf{d}_{k+1}$ satisfying
\begin{displaymath}
  \mathsf{a}(x)\mult\mathsf{d}(x) = \sum_{i=1}^{k+1} 2^{16(i-1)}\mult\mathsf{d}_i(x), \qquad x \in \BoolHCube{\mu},
\end{displaymath}
which a single \cref{piop:zerocheck} enforces; the activator collapses inactive rows to zero, so the identity pins the limb decomposition on active rows. A \cref{piop:lookup-check} of each limb into the appropriate small index column then forces $\mathsf{d}(x) \in [0, 2^n - 1]$ at every active $x$.
\begin{DefinePIOP}{Range Check}{range-check}
  \emph{Inputs:} $\Oracle{\MLE{\col}} = (\Oracle{\MLE{\mathsf{a}}}, \Oracle{\MLE{\mathsf{d}}})$, width $n$.
  \begin{enumerate}[nosep]
    \item $\PIOPProver$ and $\PIOPVerifier$ compute $\beta, k$ such that $16k + \beta = n$ and $\beta \leq 16$.
    \item $\PIOPProver$ constructs and sends the limb oracles $\Oracle{\MLE{\mathsf{d}_1}}, \ldots, \Oracle{\MLE{\mathsf{d}_{k+1}}}$, the $16$-bit limb decomposition of $\mathsf{d}$.
    \item $\PIOPProver$ and $\PIOPVerifier$ invoke \cref{piop:zerocheck} on the column $\mathsf{a}\mult\mathsf{d} - \sum_{i=1}^{k+1} 2^{16(i-1)}\mult\mathsf{d}_i$.
    \item For each $i \in [k]$, $\PIOPProver$ and $\PIOPVerifier$ invoke \cref{piop:lookup-check} on $(\mathbf{1}, \mathsf{d}_i)$ and $\mathsf{rng}_{16} := (\mathbf{1}, \idx_{16})$, proving each limb lies in $[0, 2^{16}-1]$.
    \item $\PIOPProver$ and $\PIOPVerifier$ invoke \cref{piop:lookup-check} on $(\mathbf{1}, \mathsf{d}_{k+1})$ and $\mathsf{rng}_\beta := (\mathbf{1}, \idx_\beta)$, proving the residual limb lies in $[0, 2^\beta-1]$.
  \end{enumerate}
\end{DefinePIOP}

\subsection{Sign Check}
\label{sec:sign-check}
The \emph{Sign Check} takes a column $\col = (\mathsf{d}, \mathsf{a})$, a width $\alpha$, and a sign mode $\sigma \in \{\mathsf{NonNeg}, \mathsf{Pos}, \mathsf{NonPos}, \mathsf{Neg}\}$, and proves that every active entry of $\col$ lies in the width-$\alpha$ slice of the corresponding sign. Each mode is determined by two parameters: a \emph{direction} $s_\sigma \in \{+1, -1\}$, which is $+1$ for the positive-side modes $\mathsf{NonNeg}, \mathsf{Pos}$ and $-1$ for the negative-side modes $\mathsf{NonPos}, \mathsf{Neg}$; and a lower bound $\ell_\sigma \in \{0, 1\}$, which is $0$ for the non-strict modes $\mathsf{NonNeg}, \mathsf{NonPos}$ and $1$ for the strict modes $\mathsf{Pos}, \mathsf{Neg}$ that exclude zero. Concretely,
\begin{center}\small
\begin{tabular}{lcccc}
\toprule
$\sigma$ & $\mathsf{NonNeg}$ & $\mathsf{Pos}$ & $\mathsf{NonPos}$ & $\mathsf{Neg}$ \\
\midrule
certified range & $[0, 2^\alpha-1]$ & $[1, 2^\alpha-1]$ & $[-(2^\alpha-1), 0]$ & $[-(2^\alpha-1), -1]$ \\
$(s_\sigma, \ell_\sigma)$ & $(+1, 0)$ & $(+1, 1)$ & $(-1, 0)$ & $(-1, 1)$ \\
\bottomrule
\end{tabular}
\end{center}
where negative values are taken in $\F$ via the additive inverse. In every mode the certified condition is $s_\sigma\mult\mathsf{d}[i] \in [\ell_\sigma, 2^\alpha-1]$, so we package the check as the relation
\begin{displaymath}
\NPRelation_{\mathsf{Sign}} = \left\{\,
  \NPInstance = (n, \alpha, \sigma),\;
  \NPInstancey = (\Oracle{\MLE{\col}}),\;
  \NPWitness = (\col)
  \;\middle|\;
    \forall i,\, \mathsf{a}[i] = 1 \implies \ell_\sigma \leq s_\sigma\mult\mathsf{d}[i] \leq 2^\alpha - 1
\,\right\}.
\end{displaymath}
We assume $\alpha < \log_2 |\F|$, so the range embeds unambiguously into $\F$.

The whole check reduces to a non-negativity check on the \emph{magnitude} column $\mathsf{m} := s_\sigma\mult\mathsf{d}$, which the verifier obtains by negating the oracle $\MLE{\mathsf{d}}$ in the negative-side modes. The index polynomial $\idx_\alpha$, evaluated over $\BoolHCube{\alpha}$, takes exactly the values $\{0, 1, \ldots, 2^\alpha - 1\}$. Treating it as the data of an $\alpha$-variable column with the all-ones activator gives a range column $\mathsf{rng}_\alpha := (\mathbf{1}, \idx_\alpha)$ whose value multiset is precisely $[0, 2^\alpha - 1]$, so a single \cref{piop:lookup-check} of $\mathsf{m}$ into $\mathsf{rng}_\alpha$ is already an exact range check. This naive lookup is wasteful when $\alpha$ is large, since $\mathsf{rng}_\alpha$ has $2^\alpha$ rows. We do better by splitting each entry into $16$-bit \emph{limbs} and range-checking each limb against a small index column. Choosing $k$ and $\beta$ with $\alpha = 16k + \beta$ and $\beta \leq 16$, the prover sends limb columns $\mathsf{m}_1, \ldots, \mathsf{m}_{k+1}$ satisfying
\begin{displaymath}
  \mathsf{a}(x) \mult s_\sigma \mult \mathsf{d}(x) = \sum_{i=1}^{k+1} 2^{16(i-1)} \mult \mathsf{m}_i(x), \qquad x \in \BoolHCube{\mu},
\end{displaymath}
which a single \cref{piop:zerocheck} enforces; the activator on the left collapses inactive rows to zero, so the identity pins down the limb decomposition on active rows and forces the limbs to combine to zero on inactive rows. Range-checking each limb then pins $s_\sigma\mult\mathsf{d}(x) \in [0, 2^\alpha - 1]$ at every active $x$. In the strict modes $\mathsf{Pos}$ and $\mathsf{Neg}$, a final no-zero check rules out the value $0$, narrowing the certified range to $[1, 2^\alpha - 1]$.
\begin{DefinePIOP}{Sign Check}{sign-check}
  \emph{Inputs:} $\Oracle{\MLE{\col}} = (\Oracle{\MLE{\mathsf{a}}}, \Oracle{\MLE{\mathsf{d}}})$, width $\alpha$, sign mode $\sigma \in \{\mathsf{NonNeg}, \mathsf{Pos}, \mathsf{NonPos}, \mathsf{Neg}\}$.
  \begin{enumerate}[nosep]
    \item $\PIOPProver$ and $\PIOPVerifier$ read off $s_\sigma \in \{+1,-1\}$ and $\ell_\sigma \in \{0,1\}$ from $\sigma$, and compute $\beta, k$ such that $16k + \beta = \alpha$ and $\beta \leq 16$.
    \item $\PIOPProver$ constructs and sends the limb oracles $\Oracle{\MLE{\mathsf{m}_1}}, \ldots, \Oracle{\MLE{\mathsf{m}_{k+1}}}$ of the magnitude $\mathsf{m} = s_\sigma\mult\mathsf{d}$.
    \item $\PIOPProver$ and $\PIOPVerifier$ invoke \cref{piop:zerocheck} to prove $\mathsf{a}\mult s_\sigma\mult\mathsf{d} - \sum_{i=1}^{k+1} 2^{16(i-1)}\mult\mathsf{m}_i = 0$.
    \item For each $i \in [k]$, $\PIOPProver$ and $\PIOPVerifier$ invoke \cref{piop:lookup-check} on $(\mathbf{1}, \mathsf{m}_i)$ and $\mathsf{rng}_{16} := (\mathbf{1}, \idx_{16})$, proving each limb lies in $[0, 2^{16}-1]$.
    \item $\PIOPProver$ and $\PIOPVerifier$ invoke \cref{piop:lookup-check} on $(\mathbf{1}, \mathsf{m}_{k+1})$ and $\mathsf{rng}_\beta := (\mathbf{1}, \idx_\beta)$, proving the residual limb lies in $[0, 2^\beta-1]$.
    \item If $\ell_\sigma = 1$ (the strict modes $\mathsf{Pos}, \mathsf{Neg}$), $\PIOPProver$ and $\PIOPVerifier$ additionally invoke \cref{piop:nozero-check} on $\col$ to rule out zero entries, narrowing the certified range to $[1, 2^\alpha - 1]$.
  \end{enumerate}
\end{DefinePIOP}

\subsection{Keyed Sumcheck}
The \emph{Keyed Sumcheck} takes two tables $\Left$ and $\Right$, each with a key column $\key$, a value column $\val$, and an activator column $\act$, and proves that for every possible key $\elem \in \F$, the sum of active values in $\Left$ with key $\elem$ equals the sum of active values in $\Right$ with key $\elem$. Define the keyed sums
\begin{displaymath}
    S_\Left(\elem) := \sum_{\substack{i \in [|\Left|] \\ \Left.\key[i] = \elem}} \Left.\act[i]\mult\Left.\val[i]
    \qquad\text{and}\qquad
    S_\Right(\elem) := \sum_{\substack{j \in [|\Right|] \\ \Right.\key[j] = \elem}} \Right.\act[j]\mult\Right.\val[j].
\end{displaymath}
We package this as the relation
\begin{displaymath}
\NPRelation_{\mathsf{ksum}} = \left\{\,
  \NPInstance = \bot,\;
  \NPInstancey = (\Oracle{\MLE{\Left}}, \Oracle{\MLE{\Right}}),\;
  \NPWitness = (\Left, \Right)
  \;\middle|\;
    \forall\, \elem \in \F : S_\Left(\elem) = S_\Right(\elem)
\,\right\}.
\end{displaymath}

\cref{piop:keyedsumcheck} reduces the per-key equality $\forall \elem : S_\Left(\elem) = S_\Right(\elem)$ to a single rational identity. The protocol is a generalization of the LogUp lookup protocol~\cite{Habock22}, and we derive its security from a core lemma about function equality that is implicit in LogUp.
\begin{ilemma}[function equality~\cite{Habock22}]
    \label{lemma:func-equality}
    For any functions $S_1, S_2: \F \to \F$, $S_1(\elem) = S_2(\elem)$ for every $\elem \in \F$ if and only if
    \begin{displaymath}
        \sum_{\elem \in \F} \frac{S_1(\elem)}{X - \elem} = \sum_{\elem \in \F} \frac{S_2(\elem)}{X - \elem}
    \end{displaymath}
    as rational functions in the formal variable $X$.
\end{ilemma}
Applying Lemma~\ref{lemma:func-equality} to $S_\Left$ and $S_\Right$ and collapsing the per-key double sums into single sums over rows (the per-key index sets partition the rows), the property holds if and only if, with $\mu$ and $\nu$ the hypercube dimensions of $\MLE{\Left}$ and $\MLE{\Right}$,
\begin{equation}
    \label{eq:keyedsumcheck_rational}
    \sum_{x \in \BoolHCube{\mu}} \frac{\MLE{\Left}.\act(x) \mult \MLE{\Left}.\val(x)}{X - \MLE{\Left}.\key(x)}
    \;=\;
    \sum_{x \in \BoolHCube{\nu}} \frac{\MLE{\Right}.\act(x) \mult \MLE{\Right}.\val(x)}{X - \MLE{\Right}.\key(x)}.
\end{equation}
We verify this rational identity by Schwartz--Zippel: the verifier samples a random $r \samples \F$, the prover sends the purported common value $s$ that both sides take at $X = r$, and the parties run \cref{piop:ratsumcheck} twice to confirm that both sides actually evaluate to $s$ at $X = r$.
\begin{DefinePIOP}{KeyedSumcheck}{keyedsumcheck}
  \emph{Inputs:} $\Oracle{\MLE{\Left}}, \Oracle{\MLE{\Right}}$.
    \begin{enumerate}[nosep]
        \item $\PIOPVerifier$ samples a random point $r \samples \F$ and sends it to $\PIOPProver$.
        \item $\PIOPProver$ computes and sends $\PIOPVerifier$ the value $s$, the purported evaluation of both sides of \cref{eq:keyedsumcheck_rational} at $X = r$.
        \item $\PIOPProver$ and $\PIOPVerifier$ invoke \cref{piop:ratsumcheck} twice to prove that both sides of \cref{eq:keyedsumcheck_rational} evaluate to $s$ at $X = r$.
    \end{enumerate}
\end{DefinePIOP}

\subsection{Multiset Equality Check}
The \emph{Multiset Equality Check} takes two columns $\col_\Left$ and $\col_\Right$ and proves that they hold the same multiset of active values. We package this as the relation
\begin{displaymath}
\NPRelation_{\mathsf{mset}} = \left\{\,
  \NPInstance = \bot,\;
  \NPInstancey = (\Oracle{\MLE{\col_\Left}}, \Oracle{\MLE{\col_\Right}}),\;
  \NPWitness = (\col_\Left, \col_\Right)
  \;\middle|\;
  \{\!\!\{\, \col_\Left[i] \,\}\!\!\}_i = \{\!\!\{\, \col_\Right[j] \,\}\!\!\}_j
\,\right\}.
\end{displaymath}
Two multisets agree exactly when every value $\elem \in \F$ occurs with the same multiplicity on both sides, i.e.\ $\forall\, \elem \in \F : \sum_{i:\,\col_\Left[i]=\elem} 1 = \sum_{j:\,\col_\Right[j]=\elem} 1$. This is precisely an $\elem$-keyed sum equality between two tables whose key is the column data and whose value is the constant $\mathbf{1}$. The parties locally form the auxiliary tables
\begin{displaymath}
    \Left := \bigl(\act = \col_\Left.\act,\; \key = \col_\Left.\mathsf{data},\; \val = \mathbf{1}\bigr)
    \quad\text{and}\quad
    \Right := \bigl(\act = \col_\Right.\act,\; \key = \col_\Right.\mathsf{data},\; \val = \mathbf{1}\bigr),
\end{displaymath}
and a single \cref{piop:keyedsumcheck} on $\Left$ and $\Right$ closes out the proof, with completeness and soundness inherited from the keyed sumcheck.
\begin{DefinePIOP}{Multiset Equality Check}{mset}
  \emph{Inputs:} $\Oracle{\MLE{\col_\Left}}, \Oracle{\MLE{\col_\Right}}$.
  \begin{enumerate}[nosep]
    \item $\PIOPProver$ and $\PIOPVerifier$ locally form the tables $\Left$ and $\Right$ as above.
    \item $\PIOPProver$ and $\PIOPVerifier$ invoke \cref{piop:keyedsumcheck} on $\Left$ and $\Right$.
  \end{enumerate}
\end{DefinePIOP}

\parhead{Handling tables} The foregoing column protocol extends to prove that two \emph{tables} hold the same multiset of \emph{rows}: the parties first fingerprint each table to a single column with a verifier-chosen random linear combination of its data columns, and then invoke the foregoing check on the two fingerprint columns.

\subsection{Rotation Check}
The \emph{Rotation Check} takes two columns $\col$ and $\col'$ together with a shift $t \in \mathbb{Z}$, and proves that $\col'$ is the cyclic rotation of $\col$ by $t$ positions, written $\col' = \rho_t(\col)$, where $\rho_t$ sends the entry at index $p$ to index $p+t$ modulo the column length $s$. We package this as the relation
\begin{displaymath}
\NPRelation_{\mathsf{rot}} = \left\{\,
  \NPInstance = t,\;
  \NPInstancey = (\Oracle{\MLE{\col}}, \Oracle{\MLE{\col'}}),\;
  \NPWitness = (\col, \col')
  \;\middle|\;
  \begin{array}{l}
    |\col| = |\col'| =: s \\[3pt]
    \forall\, q \in [s] : \col'[q] = \col[(q - t) \bmod s]
  \end{array}
\,\right\}.
\end{displaymath}
We first adopt the contiguous convention so that the active rows of both columns occupy hypercube indices $0, \ldots, s-1$, which the verifier enforces by asserting $\col.\act = \col'.\act = \contig_{\mu, s}$. We then reduce the rotation condition to a single multiset equality on \emph{position-tagged} columns: tag each entry of $\col$ with its own position, and tag each entry of $\col'$ with the position it should occupy in $\col$. Concretely, the parties locally construct the tagged tables
\begin{displaymath}
  \Table := \bigl(\col.\act,\; \col.\mathsf{data},\; \idx_\mu\bigr)
  \quad\text{and}\quad
  \Table' := \bigl(\col'.\act,\; \col'.\mathsf{data},\; \idxrot\bigr),
\end{displaymath}
where $\idx_\mu$ maps each $x \in \BoolHCube{\mu}$ to its hypercube position and $\idxrot$ is the \emph{shifted} index polynomial sending each active position $q$ to $(q - t) \bmod s$. Writing $t_0 := t \bmod s$, this shift is realized exactly by
\begin{displaymath}
  \idxrot(x) := \idx_\mu(x) - t_0 + s\mult\contig_{\mu, t_0}(x),
\end{displaymath}
since $\contig_{\mu, t_0}(x) = 1$ precisely when $\idx_\mu(x) < t_0$, supplying the $+s$ wraparound exactly on the positions that underflow. Because $q \mapsto (q - t) \bmod s$ is a bijection on $\{0, \ldots, s-1\}$, the tagged tables agree as multisets of rows if and only if $\col'[q] = \col[(q-t) \bmod s]$ for every $q$, which is exactly the rotation condition. A single \cref{piop:mset} on $\Table$ and $\Table'$ then proves the claim, with completeness and soundness inherited from the multiset equality check.
\begin{DefinePIOP}{Rotation Check}{rotcheck}
  \emph{Inputs:} $\Oracle{\MLE{\col}}, \Oracle{\MLE{\col'}}, t$.
  \begin{enumerate}[nosep]
    \item $\PIOPProver$ sends $\PIOPVerifier$ the active length $s$, and $\PIOPVerifier$ asserts $\col.\act = \col'.\act = \contig_{\mu, s}$.
    \item $\PIOPProver$ and $\PIOPVerifier$ locally construct the tagged tables $\Table$ and $\Table'$ as above, using the shifted index polynomial $\idxrot(x) = \idx_\mu(x) - t_0 + s\mult\contig_{\mu, t_0}(x)$ with $t_0 := t \bmod s$.
    \item $\PIOPProver$ and $\PIOPVerifier$ invoke \cref{piop:mset} to prove that $\Table$ and $\Table'$ are equal as multisets of rows.
  \end{enumerate}
\end{DefinePIOP}
